\scanpage{54}
% archive sheet number: 53
\mspage{2 - 5 -}\hfill n\textsuperscript{o} 392

\[
\underset{=}{a}\underset{=}{b}=\handwritten{\prod}\unsure{\underset{=}{p}|\underline{p}|p}^{m(\underset{=}{p})+n(\underset{=}{p})}\tag*{(2)}
\]
\[
\underset{=}{a}:\underset{=}{b}=\underset{=}{a}\underset{=}{b}^{-1}=\handwritten{\prod}\underset{=}{p}^{m(\underset{=}{p})-n(\underset{=}{p})}\tag*{(3)}
\]
\[
\underset{=}{a}+\underset{=}{b}=\handwritten{\prod}\underset{=}{p}^{\inf(m(\underset{=}{p}),n(\underset{=}{p}))}\tag*{(4)}
\]
\[
\underset{=}{a}\handwritten{\cap}\underset{=}{b}=\handwritten{\prod}\underset{=}{p}^{\sup(m(\underset{=}{p}),n(\underset{=}{p}))}.\tag*{(5)}
\]
En effet, la relation (2) est évidente ; la relation (3) en résulte, l'égalité $\underset{=}{a}:\underset{=}{b}=\underset{=}{a}\underset{=}{b}^{-1}$ découlant de la formule $\operatorname{div}(\underset{=}{a}:\underset{=}{b})=\operatorname{div}\underset{=}{a}-\operatorname{div}\underset{=}{b}$, (§ 1, n\textsuperscript{o} 2, cor. du th. 1) ; les formules (4) et (5) résultent de la prop. 2, § 1, n\textsuperscript{o} 1.

Ces résultats s'appliquent notamment à la clôture intégrale de $\unsure{\underset{=}{Z}|\mathbb Z}$ dans une extension de degré fini de $Q$.

Lorsque $A$ est principal, les résultats ci-dessus redonnent ceux d'\uline{Alg}., chap. VII, § 1, n\textsuperscript{o} 3.

4. \uline{Théorème d'approximation dans les anneaux de Dedekind}.

Dans les anneaux de Dedekind, on a un ``théorème d'approximation'' qui améliore à la fois le th. 1 du chap. VI, § 7, n\textsuperscript{o} 2 et la prop. 9 du § 1, n\textsuperscript{o} 5 :

PROPOSITION 2.- \uline{Soient} $A$ \uline{un anneau de Dedekind}, $K$ \uline{son corps des fractions}, \uline{et} $P$ \uline{l'ensemble des idéaux premiers non nuls de} $A$ ; \uline{pour} $\underset{=}{p}\handwritten{\in}P$, \uline{notons} $v_{\underset{=}{p}}$ \uline{la valuation essentielle correspondante de} $A$. \uline{Soient} $\underset{=}{p}_1,\ldots,\underset{=}{p}_q$ \uline{des éléments deux à deux distincts de} $P$, $n_1,\ldots,n_q$ \uline{des entiers rationnels}, \uline{et} $x_1,\ldots,x_q$ \uline{des éléments de} $K$. \uline{Il existe alors} $x\handwritten{\in}K$ \uline{tel que} $v_{\underset{=}{p}_i}(x-x_i)\handwritten{\geq}n_i$ \uline{pour} $1\handwritten{\leq}i\handwritten{\leq}q$, \uline{et que} $v_{\underset{=}{p}}(x)\handwritten{\geq}0$ \uline{pour tout} $\underset{=}{p}\handwritten{\in}P$ \uline{distinct des} $\underset{=}{p}_i$.

En remplaçant au besoin les $n_i$ par des entiers qui leurs sont supérieurs, on peut les supposer tous positifs. Examinons d'abord le cas où les $x_i$ sont dans $A$ ; tout revient évidemment à trouver un $x\handwritten{\in}A$ vérifiant les congruences
\[
x\handwritten{\equiv}x_i\ (\mathrm{mod.}\ \underset{=}{p}^{n_i})
\]
et l'existence de $x$ résulte alors du chap. II, § 1, n\textsuperscript{o} 2, prop. 5.
